% =============================================================
\chapter{บทนำ}
% =============================================================

\section{ความเป็นมาและความสำคัญของโครงงาน}

Disney Lorcana เป็นเกมการ์ดสะสมที่ Ravensburger และ The Walt Disney Company เปิดตัวในปี 2023 [1] ผู้เล่นที่ต้องการฝึกจัดเด็คหรือทดลองเล่นมักติดข้อจำกัด 3 ข้อ

\begin{enumerate}
    \item \textbf{ต้นทุนการ์ดจริงและสถานที่:} การ์ดหายากมีราคาสูงในตลาดรอง และการนัดเล่นต้องมีสถานที่ ผู้เล่นใหม่จึงทดลองเด็คได้ยาก
    \item \textbf{ต้นทุนเซิร์ฟเวอร์:} ระบบจำลองเกมที่รันบนเครื่องเดียวและเปิดตลอด 24 ชั่วโมงมีค่าใช้จ่ายคงที่แม้ไม่มีผู้ใช้ และล่มได้เมื่อมีผู้เล่นพร้อมกันจำนวนมาก
    \item \textbf{การซิงค์สถานะระหว่างผู้เล่น:} เกมการ์ดต้องแสดงตำแหน่งการ์ด สถานะ Ready/Exerted และแต้ม Lore ของทั้งสองฝ่ายให้ตรงกัน การดึงข้อมูลแบบ Polling ทำให้หน้าจอของผู้เล่นสองคนไม่ตรงกัน
\end{enumerate}

โครงงานนี้สร้างเว็บแอปที่จำลองกระดานเล่น Lorcana และวิเคราะห์เด็คบน AWS โดยใช้ทั้งบริการแบบ IaaS (VPC, ALB, EC2 Auto Scaling Group) และแบบ Serverless (API Gateway, Lambda, DynamoDB, SQS, SNS) งานแบ่งเป็น 3 ระยะ ระยะที่ 1 เสนอโครงงานและตั้งค่าเริ่มต้น ระยะที่ 2 สร้างหน้าเว็บ ระบบสมาชิก ห้องเล่น และสถาปัตยกรรม Multi-AZ ระยะที่ 3 เพิ่มการแยกส่วนประมวลผลเหตุการณ์ แคช ตัวขับเคลื่อนกฎเกม และการทดลองวัดผลบน AWS

\section{วัตถุประสงค์}

\begin{enumerate}
    \item สร้างกระดานจำลองการเล่น Disney Lorcana สำหรับผู้เล่น 2 คนบนเว็บเบราว์เซอร์ พร้อมห้องเล่นที่ซิงค์สถานะผ่าน WebSocket
    \item ออกแบบโครงสร้างพื้นฐานแบบ Multi-AZ ที่ทนต่อการล่มของเครื่องเดียวและขยายจำนวนเครื่องอัตโนมัติตาม CPU
    \item สร้างระบบจัดการเด็ค (Deck Builder) ที่ตรวจกฎของเด็คและวิเคราะห์ความสมดุลของเด็คแบบ Asynchronous ผ่าน Amazon SQS
    \item แยกการบันทึกผลการแข่งขันออกเป็นส่วนอิสระด้วย Amazon SNS, SQS และ Dead Letter Queue
    \item สร้างตัวขับเคลื่อนกฎเกมที่ทดสอบได้โดยไม่ต้องใช้เบราว์เซอร์ รองรับคีย์เวิร์ดทางการ 14 ชนิด
    \item ทดลองและเก็บหลักฐานการทำงานของระบบบน AWS Academy Learner Lab ภายใต้งบ 50 ดอลลาร์
\end{enumerate}

\section{ขอบเขตของระบบ}

\begin{enumerate}
    \item \textbf{สมาชิกและความปลอดภัย:} สมัคร เข้าสู่ระบบ ด้วยรหัสผ่านที่แฮชด้วย bcrypt และใช้ JSON Web Token
    \item \textbf{จัดการเด็ค:} สร้าง แก้ไข ลบเด็คส่วนตัว พร้อมตรวจกฎเด็ค (60 ใบ ไม่เกิน 2 สี ซ้ำไม่เกิน 4 ใบ) และเก็บใน Amazon DynamoDB
    \item \textbf{เปิดซองการ์ดจำลอง (Booster Pack):} สุ่มการ์ด 12 ใบต่อซอง
    \item \textbf{ตรวจการ์ดแบบ 3 มิติ:} ดูการ์ดพร้อมเอฟเฟกต์ฟอยล์
    \item \textbf{ห้องเล่นเรียลไทม์:} สร้างห้องด้วยรหัส 6 หลัก ทอยลูกเต๋าเลือกลำดับ สลับการ์ดเริ่มต้น (Mulligan) เล่นตามเทิร์น กลับเข้าห้องได้ภายใน 60 วินาทีหลังเชื่อมต่อหลุด
    \item \textbf{วิเคราะห์เด็ค:} คำนวณ Ink Curve และการกระจายหมึกผ่านคิว SQS และ Lambda
    \item \textbf{ประวัติการแข่งขัน สถิติผู้เล่น และตารางผู้นำ}
    \item \textbf{ตัวขับเคลื่อนกฎเกม:} ตรวจกติกาการลงการ์ด Quest Challenge Shift Location Song Item และ Ability
    \item \textbf{ตรวจสอบระบบ:} ติดตามผ่าน Amazon CloudWatch
\end{enumerate}

ขอบเขตที่ไม่ได้ทำ: ตัวขับเคลื่อนกฎทำงานฝั่งเบราว์เซอร์ ยังไม่มีเซิร์ฟเวอร์กลางตรวจสถานะเกม (ดูข้อจำกัดในบทที่ \ref{ch:limits}) ระบบไม่มี HTTPS เพราะ Learner Lab ไม่มีโดเมนสำหรับขอ certificate

\section{ประโยชน์ที่คาดว่าจะได้รับ}

\begin{enumerate}
    \item ผู้เล่นฝึกและทดสอบเด็คผ่านเบราว์เซอร์ได้โดยไม่ต้องติดตั้งโปรแกรม
    \item ทีมได้ประสบการณ์ออกแบบและ deploy ระบบ IaaS ร่วมกับ Serverless บน AWS พร้อมหลักฐานการวัดผลจริง
    \item มีบันทึกข้อจำกัดของ AWS Learner Lab และวิธีรับมือ ซึ่งนำไปใช้กับโครงงานอื่นได้
\end{enumerate}

\section{โครงสร้างของรายงาน}

\begin{itemize}
    \item \textbf{บทที่ 1 บทนำ:} ความเป็นมา วัตถุประสงค์ ขอบเขต
    \item \textbf{บทที่ 2 การบริหารโครงการและเครื่องมือ:} แผนงาน สถานะแต่ละ Sprint เครื่องมือที่ใช้ และกติกาเกมที่ระบบต้องจำลอง
    \item \textbf{บทที่ 3 การวิเคราะห์และออกแบบ:} ข้อกำหนด สถาปัตยกรรม ฐานข้อมูล วงจรชีวิตห้องเล่น ตัวขับเคลื่อนกฎ และเสาหลัก Well-Architected
    \item \textbf{บทที่ 4 การพัฒนาและผลการทดลอง:} หน้าจอระบบ การทดสอบอัตโนมัติ และการทดลอง 10 หัวข้อบน AWS
    \item \textbf{บทที่ \ref{ch:limits} ข้อจำกัดและค่าใช้จ่าย:} ข้อจำกัดที่ทราบและค่าใช้จ่ายจริง
    \item \textbf{บทที่ \ref{ch:summary} สรุปและข้อเสนอแนะ}
\end{itemize}
